% frechet_derivatives.tex — Fréchet Derivatives of the Kennett Reflectivity
% Compile with: lualatex frechet_derivatives.tex
\documentclass[11pt,a4paper]{article}

\usepackage{fontspec}
\setmainfont{Latin Modern Roman}
\setmonofont{Latin Modern Mono}

\usepackage{amsmath,amssymb,amsthm}
\usepackage{mathtools}
\usepackage{bm}
\usepackage{booktabs}
\usepackage[svgnames]{xcolor}
\usepackage{geometry}
\geometry{margin=2.5cm}

\usepackage{hyperref}
\hypersetup{
  colorlinks=true,
  linkcolor=DarkBlue,
  citecolor=DarkBlue,
  urlcolor=DarkBlue,
}

\usepackage{enumitem}
\usepackage{fancyhdr}
\pagestyle{fancy}
\fancyhf{}
\fancyhead[L]{\small Derivatives and Inversion Theory for Plane-Wave Reflectivity}
\fancyhead[R]{\small\thepage}
\renewcommand{\headrulewidth}{0.4pt}

% Notation shortcuts
\newcommand{\dR}{\delta\mathbf{R}}
\newcommand{\dT}{\delta\mathbf{T}}
\newcommand{\dE}{\delta\mathbf{E}}
\newcommand{\dM}{\delta\mathbf{M}}
\newcommand{\dU}{\delta\mathbf{U}}
\newcommand{\dZ}{\delta\mathbf{Z}}
\newcommand{\dRR}{\delta\mathbf{RR}}
\newcommand{\RR}{\mathbf{RR}}
\newcommand{\Rd}{\mathbf{R}_d}
\newcommand{\Ru}{\mathbf{R}_u}
\newcommand{\Tu}{\mathbf{T}_u}
\newcommand{\Td}{\mathbf{T}_d}
\newcommand{\vE}{\mathbf{E}}
\newcommand{\vM}{\mathbf{M}}
\newcommand{\vU}{\mathbf{U}}
\newcommand{\vZ}{\mathbf{Z}}
\newcommand{\vI}{\mathbf{I}}
\newcommand{\vW}{\mathbf{W}}
\newcommand{\dd}{\mathrm{d}}
\newcommand{\ppd}[2]{\frac{\partial #1}{\partial #2}}

\theoremstyle{plain}
\newtheorem{proposition}{Proposition}
\newtheorem{lemma}[proposition]{Lemma}
\theoremstyle{remark}
\newtheorem{remark}{Remark}

\title{%
  Derivatives and Inversion Theory for Plane-Wave Reflectivity\\[6pt]
  \large Kennett Recursion, Global Matrix Implicit Differentiation,\\
  and Levenberg--Marquardt Inversion%
}
\author{}
\date{}

\begin{document}
\maketitle
\thispagestyle{fancy}

\begin{abstract}
We present the complete derivative theory for plane-wave PP reflectivity
inversion in stratified elastic media.

\textbf{Part~I} (Sections~1--8) derives the Fréchet (first-order) derivatives
via tangent-linear differentiation of the Kennett addition formula, equivalent
to the analytical kernels of Dietrich \& Kormendi (1990).  The structure of
the second-order (Hessian) derivatives is presented, identifying the
cross-coupling terms that make same-layer Hessian elements combinatorially
complex.

\textbf{Part~II} (Section~9) derives the Jacobian and Hessian of the
Global Matrix Method (Chin, Hedstrom \& Thigpen, 1984) via implicit
differentiation of the linear system $\mathbf{G}\mathbf{x} = \mathbf{b}$,
showing that each derivative order reduces to a single back-substitution
with the already-factored $\mathbf{G}$.  The Block-Riccati sweep solver
(Section~9.5) reduces the forward solve from $O(N^3)$ to $O(N)$ while
preserving the implicit differentiation advantage through a chain of
$4\times 4$ \texttt{torch.linalg.solve} calls.

\textbf{Part~III} (Section~10) presents the Levenberg--Marquardt inversion
framework: the misfit functional, gradient, exact Hessian decomposition into
Gauss--Newton and second-order residual terms, the damped Newton update, and
the log-parameter transformation for positivity.

\medskip\noindent
\textbf{Convention:}  $\exp(-i\omega t)$ inverse Fourier transform; depth
positive downward; $\operatorname{Im}(\eta) > 0$ for vertical slowness
(radiation condition).
\end{abstract}

\tableofcontents

%=============================================================================
\section{The Kennett Reflectivity}
%=============================================================================

The PP reflectivity of a stratified elastic half-space with $N$ sub-ocean
layers is
\begin{equation}\label{eq:R}
  R(\omega, p)
  = e_{a,0}^{2} \;\RR_d^{(0)}[0,0],
\end{equation}
where $e_{a,0}^{2} = \exp(2i\omega\eta_0 h_0)$ is the two-way ocean phase and
$\RR_d^{(0)}$ is the cumulative downgoing reflection matrix at the ocean
bottom, computed by the upward recursion described below.

%=============================================================================
\section{Kennett Addition Formula}\label{sec:addition}
%=============================================================================

At interface~$i$ (between layers~$i$ and $i{+}1$), the cumulative reflection
from below is updated by
\begin{equation}\label{eq:addition}
  \RR_d^{(i)}
  = \Rd^{(i)} + \Tu^{(i)}\,\vM\,\vU\,\Td^{(i)},
\end{equation}
where
\begin{equation}\label{eq:MU}
  \vM = \vE_{i+1}\,\RR_d^{(i+1)}\,\vE_{i+1},
  \qquad
  \vU = \bigl(\vI - \Ru^{(i)}\,\vM\bigr)^{-1},
\end{equation}
and the phase matrix is
\begin{equation}\label{eq:E}
  \vE_j = \operatorname{diag}\bigl(e_{a,j},\; e_{b,j}\bigr),
  \qquad
  e_{a,j} = \exp(i\omega\eta_j h_j),
  \quad
  e_{b,j} = \exp(i\omega\nu_j h_j).
\end{equation}
Here $\eta_j$ and $\nu_j$ are the P- and S-wave vertical slownesses,
$h_j$ is the layer thickness, and $\Rd^{(i)}$, $\Ru^{(i)}$, $\Tu^{(i)}$,
$\Td^{(i)}$ are the $2\times 2$ scattering matrices at interface~$i$.

\medskip\noindent
\textbf{Initialisation:}
$\RR_d^{(N-1)} = \mathbf{0}$ (radiation condition at the half-space).

%=============================================================================
\section{Tangent-Linear Kennett Recursion}\label{sec:tangent}
%=============================================================================

Differentiating~\eqref{eq:addition} with respect to an arbitrary model
parameter $m_j$ yields the tangent-linear recursion for $\dRR_d$:
\begin{equation}\label{eq:tl-main}
  \boxed{
  \dRR_d^{(i)}
  = \dR_d
  + \dT_u\,\vZ\,\Td
  + \Tu\,\dZ\,\Td
  + \Tu\,\vZ\,\dT_d
  }
\end{equation}
where $\vZ = \vM\,\vU$ and the subsidiary tangents are:
\begin{align}
  \dZ   &= \dM\,\vU + \vM\,\dU, \label{eq:dZ}\\[4pt]
  \dU   &= \vU\bigl(\dR_u\,\vM + \Ru\,\dM\bigr)\vU,
           \label{eq:dU}\\[4pt]
  \dM   &= \dE\,\RR_d\,\vE
         + \vE\,\dRR_d\,\vE
         + \vE\,\RR_d\,\dE.
           \label{eq:dM}
\end{align}

\begin{remark}
Equation~\eqref{eq:dU} follows from the identity
$\dd(\vW^{-1}) = \vW^{-1}\,\dd\vW\,\vW^{-1}$
applied to $\vW = \vI - \Ru\,\vM$.
\end{remark}

\noindent\textbf{Initialisation:}
$\dRR_d^{(N-1)} = \mathbf{0}$.

%=============================================================================
\section{Primitive Derivative Chains}\label{sec:primitives}
%=============================================================================

For a perturbation of parameter $m_j$ in layer~$j$, the following chains
provide the tangent inputs to the scattering matrices and phase factors.

\subsection{Complex slowness from velocity}

From the Futterman attenuation model:
\begin{equation}\label{eq:slowness}
  s = \frac{2Q^2 + 2Qi}{(1+4Q^2)\,v},
\end{equation}
where $Q$ is the quality factor (held fixed).  Since $s$ is proportional to
$1/v$,
\begin{equation}\label{eq:dsdv}
  \boxed{\ppd{s}{v} = -\frac{s}{v}.}
\end{equation}

\subsection{Vertical slowness}

With $\eta = \sqrt{s^2 - p^2}$ and the branch chosen so that
$\operatorname{Im}(\eta) > 0$:
\begin{equation}\label{eq:detads}
  \ppd{\eta}{s} = \frac{s}{\eta}.
\end{equation}
This holds for either sign of $\eta$: the branch flip $\eta\to -\eta$ leaves
$s/\eta$ invariant.

\subsection{P-wave velocity $\alpha_j$}

\begin{align}
  \ppd{s_{p,j}}{\alpha_j} &= -\frac{s_{p,j}}{\alpha_j},
  \label{eq:dsp}\\[4pt]
  \ppd{\eta_j}{\alpha_j}
    &= \frac{s_{p,j}}{\eta_j}\cdot\ppd{s_{p,j}}{\alpha_j}
     = -\frac{s_{p,j}^2}{\eta_j\,\alpha_j},
  \label{eq:deta-alpha}\\[4pt]
  \ppd{e_{a,j}}{\alpha_j}
    &= i\omega h_j\;\ppd{\eta_j}{\alpha_j}\;e_{a,j}.
  \label{eq:dea-alpha}
\end{align}

\subsection{S-wave velocity $\beta_j$}

\begin{align}
  \ppd{s_{s,j}}{\beta_j} &= -\frac{s_{s,j}}{\beta_j},
  \qquad
  \ppd{\nu_j}{\beta_j} = -\frac{s_{s,j}^2}{\nu_j\,\beta_j},
  \qquad
  \ppd{\tilde\beta_j}{\beta_j} = \frac{\tilde\beta_j}{\beta_j},
  \label{eq:dbeta}\\[4pt]
  \ppd{e_{b,j}}{\beta_j}
    &= i\omega h_j\;\ppd{\nu_j}{\beta_j}\;e_{b,j},
  \label{eq:deb-beta}
\end{align}
where $\tilde\beta_j = 1/s_{s,j}$ is the complex S-velocity entering the
scattering matrices.

\subsection{Density $\rho_j$}

Density enters the scattering matrices directly.  It does not affect the phase
factors ($\delta e_{a,j} = \delta e_{b,j} = 0$).

\subsection{Thickness $h_j$}

\begin{equation}\label{eq:dh}
  \ppd{e_{a,j}}{h_j} = i\omega\eta_j\,e_{a,j},
  \qquad
  \ppd{e_{b,j}}{h_j} = i\omega\nu_j\,e_{b,j}.
\end{equation}
Thickness does not affect the scattering matrices.

%=============================================================================
\section{Scattering Matrix Derivatives (Solid--Solid)}\label{sec:scatmat}
%=============================================================================

\subsection{Kennett intermediate variables}

At interface~$i$ between layers with properties $(\eta_1,\nu_1,\rho_1,
\tilde\beta_1)$ above and $(\eta_2,\nu_2,\rho_2,\tilde\beta_2)$ below,
define:
\begin{align}
  \mu_k &= \rho_k\,\tilde\beta_k^2, &
  \Delta\mu &= \mu_2 - \mu_1, &
  d &= 2\Delta\mu, \label{eq:mud}\\[4pt]
  a &= \Delta\rho - p^2 d, &
  b &= \rho_2 - p^2 d, &
  c &= \rho_1 + p^2 d, \label{eq:abc}\\[4pt]
  E &= b\eta_1 + c\eta_2, &
  F &= b\nu_1 + c\nu_2, \label{eq:EF}\\[4pt]
  G &= a - d\eta_1\nu_2, &
  H &= a - d\eta_2\nu_1, \label{eq:GH}\\[4pt]
  D &= EF + GHp^2. \label{eq:D}
\end{align}

\subsection{Unnormalised scattering elements}

Define normalised variables $\hat{E} = E/D$, etc.  The six independent
scattering elements are:
\begin{align}
  Q &= (b\eta_1 - c\eta_2)\,\hat{F},
  & R &= (a + d\eta_1\nu_2)\,\hat{H}\,p^2,
  \label{eq:QR}\\[4pt]
  S &= \frac{(ab + cd\eta_2\nu_2)\,p}{D},
  & T &= (b\nu_1 - c\nu_2)\,\hat{E},
  \label{eq:ST}\\[4pt]
  U &= (a + d\eta_2\nu_1)\,\hat{G}\,p^2,
  & V &= \frac{(ac + bd\eta_1\nu_1)\,p}{D}.
  \label{eq:UV}
\end{align}

\subsection{Tangent of the intermediate variables}\label{sec:tl-intermediate}

For a perturbation $(\delta\eta_1, \delta\nu_1, \delta\rho_1,
\delta\tilde\beta_1, \delta\eta_2, \delta\nu_2, \delta\rho_2,
\delta\tilde\beta_2)$:
\begin{align}
  \delta\mu_k &= \delta\rho_k\,\tilde\beta_k^2
               + 2\rho_k\tilde\beta_k\,\delta\tilde\beta_k,
  \label{eq:dmu}\\[4pt]
  \delta d &= 2(\delta\mu_2 - \delta\mu_1),
  \label{eq:dd}\\[4pt]
  \delta a &= \delta(\Delta\rho) - p^2\,\delta d,
  \label{eq:da}\\
  \delta b &= \delta\rho_2 - p^2\,\delta d,
  \label{eq:db}\\
  \delta c &= \delta\rho_1 + p^2\,\delta d,
  \label{eq:dc}\\[4pt]
  \delta E &= \delta b\,\eta_1 + b\,\delta\eta_1
            + \delta c\,\eta_2 + c\,\delta\eta_2,
  \label{eq:dE}\\
  \delta F &= \delta b\,\nu_1 + b\,\delta\nu_1
            + \delta c\,\nu_2 + c\,\delta\nu_2,
  \label{eq:dF}\\[4pt]
  \delta G &= \delta a - \delta d\,\eta_1\nu_2
            - d\,\delta\eta_1\,\nu_2 - d\,\eta_1\,\delta\nu_2,
  \label{eq:dG}\\
  \delta H &= \delta a - \delta d\,\eta_2\nu_1
            - d\,\delta\eta_2\,\nu_1 - d\,\eta_2\,\delta\nu_1,
  \label{eq:dH}\\[4pt]
  \delta D &= \delta E\,F + E\,\delta F
            + (\delta G\,H + G\,\delta H)\,p^2.
  \label{eq:dD}
\end{align}

\subsection{Tangent of the normalised variables}

Using the quotient rule $\delta(X/D) = (\delta X\cdot D - X\cdot\delta D)/D^2$:
\begin{equation}\label{eq:dhat}
  \delta\hat{E} = \frac{\delta E\cdot D - E\cdot\delta D}{D^2},
  \quad
  \delta\hat{F} = \frac{\delta F\cdot D - F\cdot\delta D}{D^2},
  \quad\text{etc.}
\end{equation}

\subsection{Tangent of the scattering elements}\label{sec:tl-scat-elem}

\begin{align}
  \delta Q &= \delta(b\eta_1 - c\eta_2)\,\hat{F}
            + (b\eta_1 - c\eta_2)\,\delta\hat{F},
  \label{eq:dQ}\\[4pt]
  \delta R &= \delta(a + d\eta_1\nu_2)\,\hat{H}\,p^2
            + (a + d\eta_1\nu_2)\,\delta\hat{H}\,p^2,
  \label{eq:dR-scat}\\[4pt]
  \delta S &= \frac{%
    \delta(ab + cd\eta_2\nu_2)\cdot D
    - (ab + cd\eta_2\nu_2)\cdot\delta D
  }{D^2}\,p,
  \label{eq:dS}\\[4pt]
  \delta T &= \delta(b\nu_1 - c\nu_2)\,\hat{E}
            + (b\nu_1 - c\nu_2)\,\delta\hat{E},
  \label{eq:dT-scat}\\[4pt]
  \delta U &= \delta(a + d\eta_2\nu_1)\,\hat{G}\,p^2
            + (a + d\eta_2\nu_1)\,\delta\hat{G}\,p^2,
  \label{eq:dU-scat}\\[4pt]
  \delta V &= \frac{%
    \delta(ac + bd\eta_1\nu_1)\cdot D
    - (ac + bd\eta_1\nu_1)\cdot\delta D
  }{D^2}\,p.
  \label{eq:dV}
\end{align}

The product-rule expansions needed in~\eqref{eq:dQ}--\eqref{eq:dV} are:
\begin{align}
  \delta(b\eta_1 - c\eta_2) &= \delta b\,\eta_1 + b\,\delta\eta_1
                              - \delta c\,\eta_2 - c\,\delta\eta_2,
  \label{eq:d-bec}\\
  \delta(a + d\eta_1\nu_2) &= \delta a + \delta d\,\eta_1\nu_2
                             + d\,\delta\eta_1\,\nu_2 + d\,\eta_1\,\delta\nu_2,
  \label{eq:d-adn}\\
  \delta(b\nu_1 - c\nu_2) &= \delta b\,\nu_1 + b\,\delta\nu_1
                            - \delta c\,\nu_2 - c\,\delta\nu_2,
  \label{eq:d-bnc}\\
  \delta(a + d\eta_2\nu_1) &= \delta a + \delta d\,\eta_2\nu_1
                             + d\,\delta\eta_2\,\nu_1 + d\,\eta_2\,\delta\nu_1,
  \label{eq:d-adn2}\\
  \delta(ab + cd\eta_2\nu_2) &= \delta a\,b + a\,\delta b
    + \delta c\,d\eta_2\nu_2 + c\,\delta d\,\eta_2\nu_2
    + cd\,\delta\eta_2\,\nu_2 + cd\,\eta_2\,\delta\nu_2,
  \label{eq:d-abcd}\\
  \delta(ac + bd\eta_1\nu_1) &= \delta a\,c + a\,\delta c
    + \delta b\,d\eta_1\nu_1 + b\,\delta d\,\eta_1\nu_1
    + bd\,\delta\eta_1\,\nu_1 + bd\,\eta_1\,\delta\nu_1.
  \label{eq:d-acbd}
\end{align}

\subsection{Tangent of the $2\times 2$ scattering matrices}\label{sec:tl-scatmat}

Define the auxiliary square-root quantities:
\begin{equation}\label{eq:sqrts}
  \hat\eta_k = \sqrt{\eta_k},\quad
  \hat\nu_k = \sqrt{\nu_k},\quad
  \hat\rho_k = \sqrt{\rho_k},\quad
  \hat{z}_{a,k} = \hat\eta_k\hat\rho_k,\quad
  \hat{z}_{b,k} = \hat\nu_k\hat\rho_k.
\end{equation}
Their tangents are $\delta\hat\eta_k = \delta\eta_k/(2\hat\eta_k)$,\;
$\delta\hat\rho_k = \delta\rho_k/(2\hat\rho_k)$, etc., and
\begin{equation}\label{eq:dz}
  \delta\hat{z}_{a,k}
  = \delta\hat\eta_k\,\hat\rho_k + \hat\eta_k\,\delta\hat\rho_k,
  \qquad
  \delta\hat{z}_{b,k}
  = \delta\hat\nu_k\,\hat\rho_k + \hat\nu_k\,\delta\hat\rho_k.
\end{equation}

The scattering matrices and their tangents are:
\begin{equation}\label{eq:dRd}
  \dR_d = \begin{pmatrix}
    \delta Q - \delta R
    & -2i\bigl(\delta\hat\eta_1\,\hat\nu_1 + \hat\eta_1\,\delta\hat\nu_1\bigr)S
      - 2i\hat\eta_1\hat\nu_1\,\delta S \\[4pt]
    \text{(sym)}
    & \delta T - \delta U
  \end{pmatrix},
\end{equation}

\begin{equation}\label{eq:dTd}
  \dT_d = \begin{pmatrix}
    2\bigl(\delta\hat{z}_{a,1}\hat{z}_{a,2} + \hat{z}_{a,1}\delta\hat{z}_{a,2}\bigr)\hat{F}
    + 2\hat{z}_{a,1}\hat{z}_{a,2}\,\delta\hat{F}
    & \cdots \\[4pt]
    \cdots
    & 2\bigl(\delta\hat{z}_{b,1}\hat{z}_{b,2} + \hat{z}_{b,1}\delta\hat{z}_{b,2}\bigr)\hat{E}
    + 2\hat{z}_{b,1}\hat{z}_{b,2}\,\delta\hat{E}
  \end{pmatrix},
\end{equation}

\begin{equation}\label{eq:dTu}
  \dT_u = (\dT_d)^T \qquad\text{(reciprocity)},
\end{equation}

\begin{equation}\label{eq:dRu}
  \dR_u = \begin{pmatrix}
    -(\delta Q + \delta U)
    & -2i\bigl(\delta\hat\eta_2\hat\nu_2 + \hat\eta_2\delta\hat\nu_2\bigr)V
      - 2i\hat\eta_2\hat\nu_2\,\delta V \\[4pt]
    \text{(sym)}
    & -(\delta T + \delta R)
  \end{pmatrix}.
\end{equation}

%=============================================================================
\section{Locality of the Perturbation}\label{sec:locality}
%=============================================================================

For a parameter $m_j$ belonging to layer~$j$, non-zero tangent inputs occur
only at:

\begin{center}
\begin{tabular}{lll}
  \toprule
  \textbf{Source} & \textbf{Affected interface} & \textbf{Role of layer~$j$}\\
  \midrule
  $\dR_d,\,\dR_u,\,\dT_u,\,\dT_d$ & $j{-}1$ & ``layer 2'' (below)\\
  $\dR_d,\,\dR_u,\,\dT_u,\,\dT_d$ & $j$ (if $< N{-}1$) & ``layer 1'' (above)\\
  $\dE_j$ & phase through layer~$j$ & $\delta e_{a,j},\,\delta e_{b,j}$\\
  \bottomrule
\end{tabular}
\end{center}

\noindent All other interfaces have
$\dR_d = \dR_u = \dT_u = \dT_d = \mathbf{0}$ and $\dE = \mathbf{0}$.
The tangent-linear recursion~\eqref{eq:tl-main} merely propagates $\dRR_d$
through the unperturbed Kennett addition formula at those interfaces.

%=============================================================================
\section{Complete Jacobian}\label{sec:jacobian}
%=============================================================================

The full Jacobian has shape $(\text{nfreq}) \times (4N{-}1)$ where $N$ is the
number of sub-ocean layers:
\begin{equation}\label{eq:Jij}
  J_{ij}
  = \ppd{R(\omega_i)}{m_j}
  = e_{a,0}^2(\omega_i) \cdot \dRR_d^{(0)}[0,0]\bigg|_{m_j}.
\end{equation}

The parameter vector is:
\begin{equation}\label{eq:params}
  \mathbf{m}
  = \bigl[\alpha_1,\dotsc,\alpha_N,\;
          \beta_1,\dotsc,\beta_N,\;
          \rho_1,\dotsc,\rho_N,\;
          h_1,\dotsc,h_{N-1}\bigr].
\end{equation}

\begin{remark}
The half-space layer ($h_N = \infty$) is excluded from the parameter vector
because its thickness is not a free parameter.
\end{remark}

%=============================================================================
\section{Equivalence to Dietrich \& Kormendi (1990)}\label{sec:DK}
%=============================================================================

The tangent-linear recursion derived above is the \emph{forward-mode automatic
differentiation} of the Kennett recursion.  It computes the exact same quantity
as the Dietrich \& Kormendi (1990) analytical Fréchet derivative, expressed as
a recursion over interfaces rather than an integral over Green's functions.

D\&K's Fréchet kernel for a continuously varying medium is
\begin{equation}\label{eq:DK-integral}
  \delta R = \int_0^\infty G^{\uparrow}(z)\;\delta\mathbf{P}(z)\;
  G^{\downarrow}(z)\,\dd z,
\end{equation}
where $G^{\uparrow}$ and $G^{\downarrow}$ are the upgoing and downgoing
Green's function operators and $\delta\mathbf{P}$ is the first-order
perturbation of the propagator.  For a stack of homogeneous layers this
integral collapses to a finite sum over the interfaces where elastic
properties are discontinuous.

\subsection{The four partitions}\label{sec:four-partitions}

At each interface~$i$ where scattering is perturbed, the tangent-linear
recursion~\eqref{eq:tl-main} decomposes the contribution to
$\dRR_d^{(i)}$ into exactly four partitions:
\begin{equation}\label{eq:4part}
  \dRR_d^{(i)}
  = \underbrace{\dR_d^{(i)}}_{P_1}
  + \underbrace{\dT_u^{(i)}\,\vZ^{(i)}\,\Td^{(i)}}_{P_2}
  + \underbrace{\Tu^{(i)}\,\dZ^{(i)}\,\Td^{(i)}}_{P_3}
  + \underbrace{\Tu^{(i)}\,\vZ^{(i)}\,\dT_d^{(i)}}_{P_4},
\end{equation}
where $\vZ^{(i)} = \vM^{(i)}\,\vU^{(i)}$ as before.  Each partition
corresponds to a distinct physical scattering pathway:

\begin{center}
\renewcommand{\arraystretch}{1.3}
\begin{tabular}{clp{7.5cm}}
  \toprule
  & \textbf{Partition} & \textbf{Physical pathway} \\
  \midrule
  $P_1$ & $\dR_d$
    & Downgoing wave reflects off the \emph{perturbed} interface.
      No interaction with the stack below. \\
  $P_2$ & $\dT_u\,\vZ\,\Td$
    & Wave transmits down (unperturbed), reverberates in the
      stack below (unperturbed), then transmits up through the
      \emph{perturbed} interface. \\
  $P_3$ & $\Tu\,\dZ\,\Td$
    & Wave enters and exits the interface unperturbed, but the
      \emph{reverberation} between $\Ru$ and the stack below is
      perturbed (contains $\dR_u$, $\dE$, and propagated
      $\dRR_d$; see below). \\
  $P_4$ & $\Tu\,\vZ\,\dT_d$
    & Wave transmits down through the \emph{perturbed} interface,
      reverberates unperturbed below, and transmits up
      unperturbed. \\
  \bottomrule
\end{tabular}
\end{center}

\subsection{Internal structure of $P_3$}\label{sec:P3}

Partition~$P_3$ is the richest because $\dZ = \dM\,\vU + \vM\,\dU$
contains three physically distinct sub-contributions via $\dM$ and
$\dU$:

\begin{enumerate}[label=(\alph*),nosep]
  \item \textbf{Phase perturbation} ($\dE$ terms in $\dM$):
    $\dE\,\RR_d\,\vE + \vE\,\RR_d\,\dE$.
    A wave traversing layer~$j$ accumulates a perturbed phase; the
    cumulative reflection from below is unperturbed.
  \item \textbf{Propagated tangent from below} ($\dRR_d$ term in $\dM$):
    $\vE\,\dRR_d\,\vE$.
    First-order perturbations originating at deeper interfaces propagate
    upward through the unperturbed phase and enter the current interface.
  \item \textbf{Upgoing reflection perturbation} ($\dR_u$ in $\dU$):
    $\vU\,\dR_u\,\vM\,\vU$.
    The reverberative coupling between the interface and the stack below
    is perturbed because $\Ru$ changes.
\end{enumerate}

\noindent Combining these with the definition $\dU =
\vU(\dR_u\,\vM + \Ru\,\dM)\vU$ from~\eqref{eq:dU}, partition~$P_3$
expands to:
\begin{equation}\label{eq:P3-expand}
  P_3 = \Tu\Bigl[
    \underbrace{\dM\,\vU}_{\text{(a)+(b)}}
    + \underbrace{\vM\,\vU\bigl(\dR_u\,\vM
      + \Ru\,\dM\bigr)\vU}_{\text{(c) + reverb of (a)+(b)}}
  \Bigr]\Td.
\end{equation}

\subsection{Correspondence to the D\&K Green's functions}

In~\eqref{eq:DK-integral}, the operators $G^{\uparrow}$ and
$G^{\downarrow}$ encode cumulative transmission through all overlying
and underlying interfaces.  In the Kennett recursive framework these
operators are not formed explicitly; instead, the upward sweep of the
tangent-linear recursion~\eqref{eq:tl-main} through \emph{unperturbed}
interfaces acts as the propagator:

\begin{center}
\renewcommand{\arraystretch}{1.3}
\begin{tabular}{lp{9cm}}
  \toprule
  \textbf{D\&K operator} & \textbf{Kennett equivalent} \\
  \midrule
  $G^{\downarrow}(z_i)$
    & Encoded in $\RR_d^{(i+1)}$: the cumulative downgoing reflection
      of the entire stack below interface~$i$. \\
  $G^{\uparrow}(z_i)$
    & Encoded in the unperturbed Kennett addition steps at
      interfaces $i{-}1, i{-}2, \dotsc, 0$ that propagate $\dRR_d$
      from interface~$i$ up to the surface. \\
  $\delta\mathbf{P}(z_i)$
    & The 4-partition decomposition: $(\dR_d,\dT_u,\dR_u,\dT_d)$ at the
      interface, plus $\dE$ in the adjacent layer.  These are the
      ``source terms'' that inject the perturbation into the recursion.\\
  \bottomrule
\end{tabular}
\end{center}

\noindent The equivalence is exact: the tangent-linear recursion computes
the same Fréchet derivative as D\&K's integral kernel, partition by
partition, with no approximation.

\begin{proposition}[Numerical validation]
The AD Jacobian (\texttt{torch.autograd}) agrees with this tangent-linear
Jacobian (NumPy) to $< 10^{-10}$ relative error, and both agree with
central finite differences to $< 10^{-5}$.
\end{proposition}

%=============================================================================
\section{Ocean-Bottom Interface (Acoustic--Elastic)}\label{sec:ocean}
%=============================================================================

The ocean-bottom interface has the same structure as solid--solid
(Section~\ref{sec:scatmat}), with the simplifications:
\begin{itemize}[nosep]
  \item $\nu_1 = 0$ (no shear in ocean),
  \item $\tilde\beta_1 = 0$,
  \item $F = b$ (not $b\nu_1 + c\nu_2$),
  \item $H = -d\eta_2$ (not $a - d\eta_2\nu_1$),
  \item $\Rd$ has only the $[0,0]$ entry nonzero (P$\to$P only from above),
  \item $\Td$ has column~1 zero (no downgoing S from ocean),
  \item $\Tu$ has row~1 zero (no upgoing S to ocean).
\end{itemize}

\noindent Ocean parameters (layer~0) are fixed:
$\delta\eta_1 = \delta\rho_1 = 0$.

%=============================================================================
\section{Second-Order Derivatives (Hessian)}\label{sec:hessian}
%=============================================================================

\subsection{Structure of the misfit Hessian}

For a scalar L2 misfit $\chi = \sum_i |R(\omega_i) - R_{\text{obs}}(\omega_i)|^2$,
the Hessian is:
\begin{equation}\label{eq:hessian}
  H_{kl}
  = \ppd{{}^2\chi}{m_k\,\partial m_l}
  = \underbrace{%
    \operatorname{Re}\!\sum_i
    \overline{\ppd{R_i}{m_k}}\,\ppd{R_i}{m_l}
  }_{\text{Gauss--Newton}}
  + \underbrace{%
    \operatorname{Re}\!\sum_i
    \overline{r_i}\,\ppd{{}^2 R_i}{m_k\,\partial m_l}
  }_{\text{second-order}},
\end{equation}
where $r_i = R(\omega_i) - R_{\text{obs}}(\omega_i)$.

The Gauss--Newton term uses only the first-order Jacobian.  The second-order
term involves $\partial^2 R/\partial m_k\,\partial m_l$, which requires
differentiating the tangent-linear recursion a second time.

\subsection{Same-layer cross-terms}\label{sec:crossterms}

When $m_k$ and $m_l$ belong to the \emph{same} layer~$j$, the second-order
tangent-linear $\delta^2\RR_d$ receives contributions from three sources:
\begin{enumerate}[nosep]
  \item[\textbf{S1}:] Scattering at interface $j{-}1$ (layer~$j$ below),
  \item[\textbf{S2}:] Scattering at interface $j$ (layer~$j$ above),
  \item[\textbf{P}:]  Phase through layer~$j$.
\end{enumerate}

The \textbf{cross-coupling terms} S1$\times$S2, S1$\times$P, and
S2$\times$P arise because the Kennett recursion chains these three
contributions together.  The key cross-term is in $\delta^2\vM$:
\begin{equation}\label{eq:d2M}
  \delta^2\vM
  = 2\,\delta_k\vE\;\delta_l\RR_d\;\vE
  + 2\,\vE\;\delta_l\RR_d\;\delta_k\vE
  + \delta_k\vE\;\RR_d\;\delta_l\vE
  + \delta_l\vE\;\RR_d\;\delta_k\vE
  + \vE\;\delta^2\RR_d\;\vE
  + \dotsb
\end{equation}
where $\delta_k$ and $\delta_l$ denote tangent-linear operations with respect
to $m_k$ and $m_l$ respectively.

The terms $\delta_k\vE\cdot\delta_l\RR_d\cdot\vE$ couple the accumulated
first-order tangent from below ($\delta_l\RR_d$) with the local phase
perturbation ($\delta_k\vE$).  This is the scattering$\times$phase cross-term
that was identified as missing from the 1991 thesis derivation.

\subsection{Magnitude of cross-terms}

Numerical experiments on a 5-layer ocean--crust model (5\% perturbation in
$\alpha_2$) show:
\begin{center}
\begin{tabular}{lc}
  \toprule
  \textbf{Cross-term} & \textbf{Relative to max diagonal}\\
  \midrule
  $\alpha\times h$ & 67\%\\
  $\alpha\times\rho$ & 17\%\\
  $\beta\times h$ & 45\%\\
  \bottomrule
\end{tabular}
\end{center}

\noindent These cross-terms are large enough that the Gauss--Newton
approximation (which neglects the second-order term entirely) and
hand-derivations that miss cross-couplings are both inadequate for
Newton-type inversion.

\subsection{Automatic differentiation resolves the problem}

The PyTorch implementation computes the exact Hessian via
\texttt{torch.autograd.functional.hessian}, which differentiates the
computational graph twice without any manual derivation.\linebreak  This automatically
captures \emph{all} cross-terms, resolving the combinatorial complexity that
prevented completion of the 1991 hand derivation.

%=============================================================================
\section{Global Matrix Method: Implicit Differentiation}\label{sec:gmm}
%=============================================================================

The Global Matrix Method (Chin, Hedstrom \& Thigpen, 1984) assembles
all interface continuity conditions into a single linear system:
\begin{equation}\label{eq:Gxb}
  \mathbf{G}(\bm{\theta})\,\mathbf{x} = \mathbf{b}(\bm{\theta}),
\end{equation}
where $\mathbf{G}$ is the global matrix encoding the displacement-stress
continuity across all interfaces, $\mathbf{b}$ is the source vector, and
$\mathbf{x}$ contains the displacement-stress coefficients at every
interface.  Both $\mathbf{G}$ and $\mathbf{b}$ depend on the model
parameter vector $\bm{\theta} = (\alpha_1,\dotsc,\alpha_N,\;
\beta_1,\dotsc,\beta_N,\; \rho_1,\dotsc,\rho_N,\;
h_1,\dotsc,h_{N-1})$.

The reflectivity is extracted from $\mathbf{x}$ by a linear projection.
Since $\mathbf{G}$ and $\mathbf{b}$ are assembled analytically from the
layer $\mathbf{E}$-matrices, all derivatives of $\mathbf{G}$ and
$\mathbf{b}$ with respect to $\bm{\theta}$ are available in closed form.
The key advantage is that derivatives of the \emph{solution}~$\mathbf{x}$
are obtained by implicit differentiation of the linear system, without
differentiating through the internal steps of the solver.

\subsection{Jacobian (first derivative)}\label{sec:gmm-jac}

Differentiate~\eqref{eq:Gxb} with respect to a model parameter $\theta_j$:
\begin{equation}\label{eq:dGxb}
  \ppd{\mathbf{G}}{\theta_j}\,\mathbf{x}
  + \mathbf{G}\,\ppd{\mathbf{x}}{\theta_j}
  = \ppd{\mathbf{b}}{\theta_j}.
\end{equation}
Rearranging:
\begin{equation}\label{eq:jac-solve}
  \boxed{
  \mathbf{G}\,\ppd{\mathbf{x}}{\theta_j}
  = \ppd{\mathbf{b}}{\theta_j}
    - \ppd{\mathbf{G}}{\theta_j}\,\mathbf{x}.
  }
\end{equation}
This is a linear system with the \emph{same coefficient matrix}
$\mathbf{G}$ as the forward problem.  Since $\mathbf{G}$ has already been
LU-factored during the forward solve, each Jacobian column costs only a
single back-substitution --- not a new factorisation.

\subsection{Hessian (second derivative)}\label{sec:gmm-hess}

Differentiate~\eqref{eq:dGxb} again with respect to a second parameter
$\theta_k$:
\begin{equation}\label{eq:d2Gxb}
  \ppd{{}^2\mathbf{G}}{\theta_j\,\partial\theta_k}\,\mathbf{x}
  + \ppd{\mathbf{G}}{\theta_j}\,\ppd{\mathbf{x}}{\theta_k}
  + \ppd{\mathbf{G}}{\theta_k}\,\ppd{\mathbf{x}}{\theta_j}
  + \mathbf{G}\,\ppd{{}^2\mathbf{x}}{\theta_j\,\partial\theta_k}
  = \ppd{{}^2\mathbf{b}}{\theta_j\,\partial\theta_k}.
\end{equation}
Rearranging:
\begin{equation}\label{eq:hess-solve}
  \boxed{
  \mathbf{G}\,\ppd{{}^2\mathbf{x}}{\theta_j\,\partial\theta_k}
  = \ppd{{}^2\mathbf{b}}{\theta_j\,\partial\theta_k}
    - \ppd{{}^2\mathbf{G}}{\theta_j\,\partial\theta_k}\,\mathbf{x}
    - \ppd{\mathbf{G}}{\theta_j}\,\ppd{\mathbf{x}}{\theta_k}
    - \ppd{\mathbf{G}}{\theta_k}\,\ppd{\mathbf{x}}{\theta_j}.
  }
\end{equation}
Again the same matrix $\mathbf{G}$ appears on the left-hand side.  The
right-hand side requires only quantities already computed:
\begin{itemize}[nosep]
  \item the forward solution $\mathbf{x}$,
  \item the first derivatives $\partial\mathbf{x}/\partial\theta_j$ and
    $\partial\mathbf{x}/\partial\theta_k$ from the Jacobian step,
  \item the first and second derivatives of $\mathbf{G}$ and $\mathbf{b}$
    with respect to $\bm{\theta}$, assembled analytically from the layer
    matrices.
\end{itemize}

\noindent For the diagonal Hessian element ($\theta_k = \theta_j$)
equation~\eqref{eq:hess-solve} simplifies to:
\begin{equation}\label{eq:hess-diag}
  \mathbf{G}\,\ppd{{}^2\mathbf{x}}{\theta_j^2}
  = \ppd{{}^2\mathbf{b}}{\theta_j^2}
    - \ppd{{}^2\mathbf{G}}{\theta_j^2}\,\mathbf{x}
    - 2\,\ppd{\mathbf{G}}{\theta_j}\,\ppd{\mathbf{x}}{\theta_j}.
\end{equation}

\subsection{Cost comparison}\label{sec:gmm-cost}

\begin{center}
\renewcommand{\arraystretch}{1.3}
\begin{tabular}{lll}
  \toprule
  \textbf{Quantity} & \textbf{Kennett recursion} & \textbf{Global Matrix} \\
  \midrule
  Forward $\mathbf{x}$
    & Recursive sweep: $O(N)$ $2\times 2$ ops
    & LU factorise $\mathbf{G}$: $O(N^3)$ \\
  Jacobian $\partial\mathbf{x}/\partial\theta_j$
    & Backprop through full graph
    & One back-substitution per $\theta_j$ \\
  Hessian $\partial^2\mathbf{x}/\partial\theta_j\partial\theta_k$
    & Backprop through the Jacobian graph
    & One back-substitution per $(\theta_j,\theta_k)$ \\
  \bottomrule
\end{tabular}
\end{center}

\noindent The Hessian is where implicit differentiation pays off most:
the Kennett method requires AD to differentiate \emph{through} the
already expensive Jacobian computation, doubling the graph depth, while
the Global Matrix Method reuses the same LU factorisation for every
derivative order.  Benchmarks show ${\sim}1.7\times$ speedup for the
Hessian, compared to ${\sim}1.2\times$ for the Jacobian
(see \texttt{GlobalMatrix/investigate\_ad.py}).

\begin{remark}
In PyTorch, \texttt{torch.linalg.solve} implements implicit
differentiation in its backward pass automatically.
\texttt{torch.func.hessian} applies the identity~\eqref{eq:hess-solve}
twice through the autograd graph without the user writing any derivative
formulae.
\end{remark}

\subsection{Block-tridiagonal structure and the Riccati connection}%
\label{sec:riccati}

The global matrix $\mathbf{G}$ is \emph{block-tridiagonal}: the
displacement-stress continuity at interface~$i$ couples only the wave
coefficients of layers~$i$ and $i{+}1$.  Partitioning $\mathbf{G}$
into $4\times 4$ blocks (P-SV up/down amplitudes per layer):
\begin{equation}\label{eq:block-tridiag}
  \mathbf{G} =
  \begin{pmatrix}
    \mathbf{A}_0 & \mathbf{C}_0 \\
    \mathbf{B}_1 & \mathbf{A}_1 & \mathbf{C}_1 \\
    & \mathbf{B}_2 & \mathbf{A}_2 & \ddots \\
    & & \ddots & \ddots & \mathbf{C}_{N-2} \\
    & & & \mathbf{B}_{N-1} & \mathbf{A}_{N-1}
  \end{pmatrix},
\end{equation}
where $\mathbf{A}_i$, $\mathbf{B}_i$, $\mathbf{C}_i$ are $4\times 4$
blocks assembled from the displacement-stress eigenvector matrices
$\mathbf{E}_i$ and phase factors of adjacent layers.

The optimal direct solver for this structure is the \textbf{block-Thomas
algorithm} (block-LU without fill-in), which sweeps upward through the
interfaces eliminating one block row at a time.  At each step, the
algorithm computes a \textbf{Schur complement}:
\begin{equation}\label{eq:schur}
  \mathbf{S}_i
  = \mathbf{A}_i - \mathbf{B}_i\,\mathbf{S}_{i+1}^{-1}\,\mathbf{C}_i,
\end{equation}
with $\mathbf{S}_{N-1} = \mathbf{A}_{N-1}$.  This is a matrix Riccati
recursion --- and it is \emph{algebraically identical} to the Kennett
addition formula~\eqref{eq:addition}.

\begin{proposition}[Kennett $\equiv$ Riccati block solver]
The Kennett recursive reflectivity method is the Riccati block solver
applied to the block-tridiagonal global matrix system.  The cumulative
reflection matrix $\RR_d^{(i)}$ is the Schur complement of the
sub-system below interface~$i$, expressed in wave-amplitude coordinates
rather than displacement-stress coordinates.  The two formulations are
related by the eigenvector transformation:
\begin{equation}\label{eq:kennett-riccati}
  \RR_d^{(i)}
  = \mathbf{P}\;\mathbf{S}_i^{-1}\;\mathbf{Q},
\end{equation}
where $\mathbf{P}$ and $\mathbf{Q}$ are projections from the
displacement-stress basis to the up/downgoing wave-amplitude basis.
\end{proposition}

This identity has an important consequence: the 1991 thesis (Nestor)
used a Riccati block solver for the forward problem, which was
\emph{already optimal} --- $O(N)$ in the number of layers, with each
step involving only $4\times 4$ block operations.  The current
implementation uses a dense $O(N^3)$ solver
(\texttt{torch.linalg.solve}), which is a step backwards for the forward
solve but gains the implicit differentiation advantage for derivatives.

\medskip\noindent
\textbf{Unifying perspective.}  The three approaches form a hierarchy:

\begin{center}
\renewcommand{\arraystretch}{1.3}
\begin{tabular}{lp{4.2cm}p{4.2cm}}
  \toprule
  \textbf{Method}
    & \textbf{Forward solve}
    & \textbf{Derivatives} \\
  \midrule
  Kennett recursion
    & $O(N)$ block-Riccati sweep
    & AD through recursive graph \\
  GMM + Riccati solver (1991 thesis)
    & $O(N)$ block-Riccati sweep
    & Hand-derived (intractable for Hessian) \\
  GMM + dense solver
    & $O(N^3)$ dense LU
    & Implicit diff via \texttt{torch.linalg.solve} \\
  GMM + Riccati solver (current)
    & $O(N)$ differentiable block-Riccati
    & Implicit diff through each $4\times 4$ block solve \\
  \bottomrule
\end{tabular}
\end{center}

\noindent This is precisely what the Riccati solver described in the
next subsection implements: the $O(N)$ forward complexity of the
block-Riccati sweep combined with PyTorch's implicit differentiation
through each $4\times 4$ \texttt{torch.linalg.solve} call.

\subsection{Differentiable Block-Riccati Solver}\label{sec:diff-riccati}

The Block-Riccati sweep reduces the global matrix solve from $O(N^3)$
to $O(N)$ per frequency while preserving differentiability via
\texttt{torch.linalg.solve}.  The key insight is that the reflectivity
can be computed by propagating a $2\times 2$ Riccati matrix through the
layer stack, where each propagation step requires solving a single
$4\times 4$ linear system.  Since \texttt{torch.linalg.solve} provides
implicit differentiation in its backward pass, the Jacobian and Hessian
are computed automatically through the chain of block solves.

\paragraph{Riccati variable.}
Define the $2\times 2$ matrix $\mathbf{Y}_k$ as the ``effective
reflection'' at the bottom of finite elastic layer~$k$: it maps the
downgoing amplitudes $\mathbf{D}_k$ (referenced at the top of layer~$k$)
to the upgoing amplitudes $\mathbf{U}_k$ (referenced at the bottom):
\begin{equation}\label{eq:riccati-var}
  \mathbf{U}_k = \mathbf{Y}_k\,\mathbf{D}_k.
\end{equation}
The Riccati matrix encodes all reflections from the layer stack below
interface~$k$.

\paragraph{Upward sweep (bottom $\to$ top).}
\textbf{Initialise:} $\mathbf{Y} = \mathbf{0}$ (radiation condition at
the half-space --- no upgoing waves from below).

\medskip\noindent
For $k = M, M{-}1, \dotsc, 1$ (where $M$ = number of finite elastic
layers), the Riccati matrix is updated by solving a $4\times 4$ linear
system at each interface:

\begin{enumerate}[nosep]
\item \textbf{Effective downgoing matrix} from the layer below ($k{+}1$):
\begin{equation}\label{eq:Eeff}
  \underset{4\times 2}{\mathbf{E}_{\text{eff}}}
  = \mathbf{E}_d^{(k+1)}
  + \mathbf{E}_u^{(k+1)}\,\operatorname{diag}(\mathbf{p}_{k+1})\,\mathbf{Y},
\end{equation}
where $\mathbf{p}_{k+1} = (e^{i\omega\eta_{k+1}h_{k+1}},\;
e^{i\omega\nu_{k+1}h_{k+1}})$ is the phase diagonal.  For the
half-space ($k{+}1 = N$): $\mathbf{E}_{\text{eff}} = \mathbf{E}_d^{(N)}$.

\item \textbf{Phased downgoing} of layer~$k$:
\begin{equation}\label{eq:Ad}
  \underset{4\times 2}{\mathbf{A}_d}
  = \mathbf{E}_d^{(k)}\,\operatorname{diag}(\mathbf{p}_k).
\end{equation}

\item \textbf{Assemble and solve} the $4\times 4$ system:
\begin{equation}\label{eq:riccati-solve}
  \boxed{
  \underbrace{\bigl[\mathbf{E}_u^{(k)} \;\big|\; {-}\mathbf{E}_{\text{eff}}\bigr]}
    _{4\times 4}
  \;\underbrace{\begin{pmatrix}\mathbf{Y}_{\text{new}} \\ \mathbf{Z}\end{pmatrix}}
    _{4\times 2}
  = -\mathbf{A}_d.
  }
\end{equation}

\item \textbf{Extract:} $\mathbf{Y} \leftarrow \mathbf{Y}_{\text{new}}$
(top two rows of the solution).
\end{enumerate}

\noindent
The physical meaning of~\eqref{eq:riccati-solve} is the continuity
condition at the bottom of layer~$k$:
\begin{equation}
  \mathbf{E}_d^{(k)}\,\operatorname{diag}(\mathbf{p}_k)\,\mathbf{D}_k
  + \mathbf{E}_u^{(k)}\,\mathbf{U}_k
  = \mathbf{E}_{\text{eff}}\,\mathbf{D}_{k+1},
\end{equation}
with the Riccati substitution $\mathbf{U}_k = \mathbf{Y}_{\text{new}}\,
\mathbf{D}_k$ and $\mathbf{D}_{k+1} = \mathbf{Z}\,\mathbf{D}_k$.

\paragraph{Ocean-bottom extraction.}
After the sweep, $\mathbf{Y} = \mathbf{Y}_1$ (the Riccati matrix at the
top of layer~1).  The effective eigenvector matrix at the ocean bottom is:
\begin{equation}\label{eq:E1eff}
  \mathbf{E}_1^{\text{eff}}
  = \mathbf{E}_d^{(1)}
  + \mathbf{E}_u^{(1)}\,\operatorname{diag}(\mathbf{p}_1)\,\mathbf{Y}_1.
\end{equation}
The ocean-bottom continuity conditions (u$_z$, $\sigma_{zz}$, and
$\sigma_{xz} = 0$) yield a $2\times 2$ system for $\mathbf{D}_1$ after
eliminating the ocean upgoing amplitude $U_0^P$:
\begin{equation}\label{eq:ocean-sys}
  \begin{pmatrix}
    e_u^{(0)}[0]\;\mathbf{E}_1^{\text{eff}}[2,:] -
    e_u^{(0)}[1]\;\mathbf{E}_1^{\text{eff}}[1,:] \\[3pt]
    \mathbf{E}_1^{\text{eff}}[3,:]
  \end{pmatrix}
  \mathbf{D}_1
  =
  \begin{pmatrix}
    \bigl(e_u^{(0)}[0]\,e_d^{(0)}[1] - e_u^{(0)}[1]\,e_d^{(0)}[0]\bigr)\,e_0 \\
    0
  \end{pmatrix},
\end{equation}
where $e_0 = \exp(i\omega\eta_0 h_0)$ is the one-way ocean phase.
The upgoing P-wave amplitude is recovered from u$_z$ continuity:
\begin{equation}\label{eq:U0P}
  U_0^P = \frac{\mathbf{E}_1^{\text{eff}}[1,:]\cdot\mathbf{D}_1
                - e_d^{(0)}[0]\,e_0}{e_u^{(0)}[0]},
\end{equation}
and the reflectivity is $R = e_0\,U_0^P$.

\paragraph{Implicit differentiation through the sweep.}
Each step of the Riccati sweep calls
\texttt{torch.linalg.solve} on a $4\times 4$
system~\eqref{eq:riccati-solve}, and the ocean extraction calls it on a
$2\times 2$ system~\eqref{eq:ocean-sys}.  PyTorch's
autograd provides implicit differentiation through each solve
automatically:

For a system $\mathbf{F}\mathbf{X} = \mathbf{C}$, the adjoint
(backward pass) computes
\begin{equation}\label{eq:implicit-adj}
  \bar{\mathbf{C}} = \mathbf{F}^{-T}\bar{\mathbf{X}},
  \qquad
  \bar{\mathbf{F}} = -\bar{\mathbf{C}}\,\mathbf{X}^T,
\end{equation}
where $\bar{(\cdot)}$ denotes the adjoint (gradient of the loss with
respect to that quantity).  The chain rule through the Riccati sweep
composes these adjoint operations backward through the interfaces,
yielding the Jacobian $\partial R/\partial\theta_j$ at cost
$O(N \times 4^3)$ per frequency --- identical complexity to the forward
sweep.

\paragraph{Hessian.}
\texttt{torch.func.hessian} computes the second-order derivatives by
applying reverse-mode AD twice: once to form the gradient of the scalar
misfit, then a second reverse pass to differentiate the gradient.  Each
level of differentiation passes through the same chain of block solves,
using the implicit differentiation identity at each step.  The total
Hessian cost is $O(N \times P^2)$ per frequency, where $P$ is the number
of model parameters --- the same as the dense solver, but with the
forward pass reduced from $O(N^3)$ to $O(N)$.

\paragraph{Complexity summary.}
\begin{center}
\renewcommand{\arraystretch}{1.3}
\begin{tabular}{lll}
  \toprule
  \textbf{Quantity}
    & \textbf{Dense GMM}
    & \textbf{Riccati GMM} \\
  \midrule
  Forward $R(\omega)$
    & $O(N_\omega \cdot N^3)$
    & $O(N_\omega \cdot N \cdot 4^3)$ \\
  Jacobian $\partial R/\partial\theta_j$
    & $O(N_\omega \cdot N^3 + N_\omega \cdot P \cdot N^2)$
    & $O(N_\omega \cdot N \cdot P \cdot 4^3)$ \\
  Hessian $\partial^2\chi/\partial\theta_j\partial\theta_k$
    & $O(N_\omega \cdot P^2 \cdot N^2)$
    & $O(N_\omega \cdot P^2 \cdot N \cdot 4^3)$ \\
  \bottomrule
\end{tabular}
\end{center}

\noindent For $N = 100$ layers, the Riccati solver is approximately
$60{,}000\times$ faster for the forward solve.  The derivative advantage
is smaller (the dense solver reuses the LU factorisation) but still
significant: $O(N)$ versus $O(N^2)$ per Jacobian column.

\paragraph{Edge cases.}
\begin{itemize}[nosep]
\item \textbf{2-layer model} ($M = 0$): The Riccati sweep is empty.
  The ocean couples directly to the half-space via
  $\mathbf{E}_1^{\text{eff}} = \mathbf{E}_d^{(1)}$.
\item \textbf{3-layer model} ($M = 1$): A single Riccati step, then
  ocean extraction.
\item \textbf{Free surface}: Applied identically after extracting $R$:
  $R_{\text{fs}} = R/(1 + R)$.
\end{itemize}

%=============================================================================
\section{Levenberg--Marquardt Inversion}\label{sec:inversion}
%=============================================================================

\subsection{Misfit functional}\label{sec:misfit}

Given observed plane-wave reflectivity $R_{\text{obs}}(\omega_i, p_s)$
at frequencies $\omega_i$ ($i = 1,\dotsc,N_\omega$) and ray parameters
$p_s$ ($s = 1,\dotsc,N_p$), the L2 misfit is:
\begin{equation}\label{eq:misfit}
  \chi(\bm{\theta})
  = \sum_{s=1}^{N_p}\sum_{i=1}^{N_\omega}
    \bigl|R(\omega_i, p_s;\bm{\theta})
          - R_{\text{obs}}(\omega_i, p_s)\bigr|^2.
\end{equation}

\subsection{Gradient}\label{sec:gradient}

The gradient of the misfit is:
\begin{equation}\label{eq:grad}
  g_k
  = \ppd{\chi}{\theta_k}
  = 2\operatorname{Re}\sum_{s,i}
    \overline{r_{si}}\;\ppd{R_{si}}{\theta_k},
\end{equation}
where $r_{si} = R(\omega_i,p_s;\bm{\theta})
- R_{\text{obs}}(\omega_i,p_s)$ is the complex residual.  In matrix
form:
\begin{equation}\label{eq:grad-mat}
  \mathbf{g} = 2\operatorname{Re}\bigl(\mathbf{J}^H\,\mathbf{r}\bigr),
\end{equation}
where $\mathbf{J}$ is the Jacobian matrix with entries
$J_{(s,i),k} = \partial R_{si}/\partial\theta_k$ and
$\mathbf{r}$ is the residual vector.

\subsection{Hessian decomposition}\label{sec:hess-decomp}

The Hessian of the misfit decomposes as (cf.~Section~\ref{sec:hessian}):
\begin{equation}\label{eq:hess-full}
  H_{kl}
  = \ppd{{}^2\chi}{\theta_k\,\partial\theta_l}
  = \underbrace{%
    2\operatorname{Re}\sum_{s,i}
    \overline{\ppd{R_{si}}{\theta_k}}\;\ppd{R_{si}}{\theta_l}
  }_{\text{Gauss--Newton: } \mathbf{H}_{\text{GN}} = 2\operatorname{Re}(\mathbf{J}^H\mathbf{J})}
  + \underbrace{%
    2\operatorname{Re}\sum_{s,i}
    \overline{r_{si}}\;\ppd{{}^2 R_{si}}{\theta_k\,\partial\theta_l}
  }_{\text{second-order residual term}}.
\end{equation}

The Gauss--Newton approximation $\mathbf{H} \approx \mathbf{H}_{\text{GN}}$
drops the second-order term.  This is valid near the solution where
$\|\mathbf{r}\| \to 0$, but fails far from the solution.

\subsection{The Levenberg--Marquardt update}\label{sec:LM}

At each iteration, the parameter update $\delta\bm{\theta}$ solves:
\begin{equation}\label{eq:LM}
  \boxed{
  \bigl(\mathbf{H} + \lambda\,\mathbf{I}\bigr)\,\delta\bm{\theta}
  = -\mathbf{g},
  }
\end{equation}
where $\mathbf{H}$ is the exact Hessian~\eqref{eq:hess-full},
$\mathbf{g}$ is the gradient~\eqref{eq:grad-mat}, and $\lambda > 0$
is the damping parameter.

\medskip\noindent
\textbf{Damping strategy:}
\begin{itemize}[nosep]
  \item If $\chi(\bm{\theta} + \delta\bm{\theta})
        < \chi(\bm{\theta})$: accept the step, decrease
        $\lambda \leftarrow \lambda/\nu$ (trust the Newton direction).
  \item If $\chi(\bm{\theta} + \delta\bm{\theta})
        \geq \chi(\bm{\theta})$: reject the step, increase
        $\lambda \leftarrow \nu\lambda$ (more regularisation).
\end{itemize}

\noindent For $\lambda \to 0$ the update approaches the full Newton step
with quadratic convergence; for $\lambda \to \infty$ it approaches the
steepest-descent direction $\delta\bm{\theta} \propto -\mathbf{g}$.

\subsection{Log-parameter transformation}\label{sec:logparams}

To enforce positivity of physical parameters and improve conditioning,
the inversion operates in log-parameter space:
\begin{equation}\label{eq:logparams}
  \bm{\phi} = \log\bm{\theta}
  = \bigl(\log\alpha_1,\dotsc,\log\alpha_N,\;
          \log\beta_1,\dotsc,\log\beta_N,\;
          \log\rho_1,\dotsc,\log\rho_N,\;
          \log h_1,\dotsc,\log h_{N-1}\bigr).
\end{equation}

The Jacobian in log-space is related to the physical-space Jacobian by:
\begin{equation}\label{eq:log-jac}
  \ppd{R}{\phi_k}
  = \ppd{R}{\theta_k}\,\theta_k
  = \theta_k\,\ppd{R}{\theta_k},
\end{equation}
which is equivalent to a diagonal scaling $\mathbf{J}_\phi
= \mathbf{J}\,\operatorname{diag}(\bm{\theta})$.  The
Hessian transforms correspondingly:
\begin{equation}\label{eq:log-hess}
  \ppd{{}^2\chi}{\phi_k\,\partial\phi_l}
  = \theta_k\,\theta_l\,\ppd{{}^2\chi}{\theta_k\,\partial\theta_l}
  + \delta_{kl}\,\theta_k\,\ppd{\chi}{\theta_k},
\end{equation}
where $\delta_{kl}$ is the Kronecker delta.  Both transformations are
handled automatically by PyTorch AD when the forward model takes
$\exp(\bm{\phi})$ as input.

\subsection{Convergence properties}\label{sec:convergence}

\begin{center}
\renewcommand{\arraystretch}{1.3}
\begin{tabular}{lcc}
  \toprule
  \textbf{Method} & \textbf{Hessian used} & \textbf{Convergence rate} \\
  \midrule
  Steepest descent & none
    & Linear: $\|\bm{\theta}_{k+1} - \bm{\theta}^*\|
      \leq c\,\|\bm{\theta}_k - \bm{\theta}^*\|$ \\
  Gauss--Newton & $\mathbf{J}^H\mathbf{J}$
    & Superlinear (near solution) \\
  Full Newton (LM) & exact $\mathbf{H}$
    & Quadratic: $\|\bm{\theta}_{k+1} - \bm{\theta}^*\|
      \leq c\,\|\bm{\theta}_k - \bm{\theta}^*\|^2$ \\
  \bottomrule
\end{tabular}
\end{center}

\noindent The 1991 thesis used the Gauss--Newton approximation because
hand-deriving the full Hessian through the Kennett recursion was
intractable.  With automatic differentiation, the exact Hessian is
obtained at no additional implementation effort --- only computational
cost --- enabling true quadratic convergence of the Levenberg--Marquardt
iteration.

%=============================================================================
\section*{References}
%=============================================================================

\begin{enumerate}[label={[\arabic*]},nosep]
  \item B.L.N.\ Kennett, \emph{Seismic Wave Propagation in Stratified Media},
    Cambridge University Press, 1983.
  \item M.\ Dietrich and F.\ Kormendi, ``Perturbation of the plane-wave
    reflectivity of a depth-dependent elastic medium by weak inhomogeneities,''
    \emph{Geophys.\ J.\ Int.}, vol.\ 100, pp.\ 203--214, 1990.
  \item F.\ Kormendi and M.\ Dietrich, ``Nonlinear waveform inversion of
    plane-wave seismograms in stratified elastic media,''
    \emph{Geophysics}, vol.\ 56, pp.\ 664--674, 1991.
  \item R.C.Y.\ Chin, G.W.\ Hedstrom and L.\ Thigpen, ``Matrix methods in
    synthetic seismograms,'' \emph{Geophys.\ J.\ R.\ astr.\ Soc.}, vol.\ 77,
    pp.\ 483--502, 1984.
  \item D.W.\ Marquardt, ``An algorithm for least-squares estimation of
    nonlinear parameters,'' \emph{SIAM J.\ Appl.\ Math.}, vol.\ 11,
    pp.\ 431--441, 1963.
  \item T.M.\ Nestor, \emph{A Practical Implementation of Waveform Inversion
    in a Stratified Marine Environment}, Master of Science thesis, Monash
    University, 1991.
\end{enumerate}

\end{document}
